\documentclass[11pt]{article}

\usepackage[T1]{fontenc}
\usepackage[utf8]{inputenc}
\usepackage[english]{babel}
\usepackage[margin=1in]{geometry}
\usepackage{lmodern}
\usepackage{microtype}
\usepackage{mathtools}
\usepackage{amsmath,amssymb,amsthm}
\usepackage{array}
\usepackage{booktabs}
\usepackage{tabularx}
\usepackage{float}
\usepackage{graphicx}
\usepackage{caption}
\usepackage{enumitem}
\usepackage{fancyhdr}
\usepackage[numbers,sort&compress]{natbib}
\usepackage{doi}
\usepackage{xurl}
\usepackage{hyperref}
\usepackage{cleveref}
\usepackage{orcidlink}
\usepackage{titling}



\DeclareMathOperator{\Capacity}{Cap}
\DeclareMathOperator{\Cost}{Cost}
\DeclareMathOperator{\Ran}{Ran}
\DeclareMathOperator{\Ker}{Ker}

\newcommand{\preceqq}{\preceq}
\newcommand{\eps}{\varepsilon}
 
\newtheorem{theorem}{Theorem}[section]
\newtheorem{lemma}[theorem]{Lemma}
\newtheorem{proposition}[theorem]{Proposition}
\newtheorem{corollary}[theorem]{Corollary}
\theoremstyle{definition}
\newtheorem{definition}[theorem]{Definition}
\newtheorem{remark}[theorem]{Remark}
\floatstyle{ruled}
\newfloat{algorithm}{tbp}{loa}
\floatname{algorithm}{Algorithm}

\urlstyle{same}
\captionsetup{
  font=small,
  labelfont=bf,
  labelsep=period
}
\captionsetup[table]{skip=6pt,position=top}
\captionsetup[figure]{skip=6pt}
\captionsetup[algorithm]{skip=6pt,position=top}
\crefname{algorithm}{algorithm}{algorithms}
\Crefname{algorithm}{Algorithm}{Algorithms}
\crefname{theorem}{theorem}{theorems}
\Crefname{theorem}{Theorem}{Theorems}
\crefname{lemma}{lemma}{lemmas}
\Crefname{lemma}{Lemma}{Lemmas}
\crefname{proposition}{proposition}{propositions}
\Crefname{proposition}{Proposition}{Propositions}
\crefname{corollary}{corollary}{corollaries}
\Crefname{corollary}{Corollary}{Corollaries}
\crefname{definition}{definition}{definitions}
\Crefname{definition}{Definition}{Definitions}
\crefname{remark}{remark}{remarks}
\Crefname{remark}{Remark}{Remarks}
\setlist[enumerate]{leftmargin=*, itemsep=2pt, topsep=4pt}
\setlength{\droptitle}{-3.2em}
\pretitle{\begin{center}\LARGE\bfseries}
\posttitle{\par\end{center}\vspace{0.5em}}
\preauthor{\begin{center}\normalsize}
\postauthor{\par\end{center}\vspace{-0.4em}}
\predate{\begin{center}\small}
\postdate{\par\end{center}\vspace{0.6em}}
\setlength{\emergencystretch}{2em}
\clubpenalty=10000
\widowpenalty=10000
\displaywidowpenalty=10000
\interfootnotelinepenalty=10000
\allowdisplaybreaks[2]
\fancypagestyle{firstpage}{\fancyhf{}
  \fancyfoot[C]{\scriptsize Automorph Inc., Wilmington, DE, USA \quad \texttt{ioannis@automorph.io} \quad \textcopyright\ 2026 \quad CC-BY 4.0}
  \renewcommand{\headrulewidth}{0pt}
  \renewcommand{\footrulewidth}{0pt}
}

\title{Adequacy Residuals and Blind-Spot Currency}
\author{Ioannis Tsiokos\,\orcidlink{0009-0009-7659-5964}}
\date{\today}
\hypersetup{
  pdftitle={Adequacy Residuals and Blind-Spot Currency},
  pdfauthor={Ioannis Tsiokos},
  hidelinks
}

\begin{document}

\maketitle
\thispagestyle{firstpage}

\begin{abstract}
Given two probe families on the same finite-dimensional carrier, this paper asks what a dissolving family \(D\) measures beyond what a native family \(L\) can explain under an audit energy \(C\).  Under the standing legal-energy quotient and null-mode legality hypotheses, the adequacy residual \(\Xi_C(D\mid L)\) is the Schur-complement leftover built from the block currencies \(K_{LL}\), \(K_{DD}\), \(K_{DL}\), and \(K_{LD}\).  We develop this residual as a calculus: the projection identity identifies it with \(T_D(I-P_L)T_D^*\) in energy-scaled coordinates; the optimal-residual and exact-adequacy theorems characterize Loewner minimality and zero residual; promotion combines native and residual budgets into a dissolving budget; and chain-rule, data-processing, and transfer laws describe how the residual changes under native extension, post-processing, coarsening, and exact or defective bridges, including the dissolving post-processing congruence.  The blind-spot witness theorem turns failure of a declared residual budget into an explicit recombination vector with positive excess.  The paper does not claim exact adequacy generically, does not classify all probe-family pairs, and does not claim novelty for the classical operator-theoretic constructions it packages; an application of the resulting calculus in a layer-dissolving emergence framework is given in \citet{Tsiokos2026NeedleKiller}.
\end{abstract}

\medskip
\noindent\textbf{Keywords.} Six Birds theory; adequacy residual; blind-spot currency; Schur complement; Moore--Penrose pseudoinverse; Loewner order; data-processing.

\section{Introduction}
\label{sec:intro}

In many finite-dimensional operator problems one has two families of
measurements on the same carrier.  The first family records what is available at
the original level of description; the second becomes available only after the
description has been enriched.  The natural question is not merely whether the
second family is large, but how much of its response is already explained by the
first family and how much remains genuinely leftover.  For a block-positive
currency matrix, the Schur complement is exactly the algebraic operation that
subtracts the explained block and leaves the unexplained one.  This paper
develops that leftover block as a first-class residual object.

A \emph{native probe family} is the linear family of measurements available at
the original layer; throughout the paper it is represented by a finite-rank map
\(L\colon E\to Y\) from a finite-dimensional carrier \(E\) to a response space
\(Y\).  A \emph{dissolving probe family} is a second linear family
\(D\colon E\to Z\), available after passing to a richer layer; the name records
its use in \citet{Tsiokos2026NeedleKiller}, where such probes are used to
dissolve lower-layer obstructions, but the present paper treats \(D\) simply as
another declared measurement family.  The carrier is priced by an audit energy
\(C=C^*\succeq 0\), zero-energy modes are handled by the legal quotient fixed in
\Cref{sec:setup}, and the currency matrices are the finite-dimensional
quadratic response operators obtained from \(C\), its Moore--Penrose
pseudoinverse, and the chosen probes.  In this language, \emph{leftover
currency} means the quantitative content of what \(D\) measures beyond what
\(L\) explains, expressed as a positive operator on the dissolving response
space \(Z\).  A \emph{blind-spot currency} is the same leftover operator viewed
as a budgeted obstruction: it is a positive residual that must be dominated by a
declared budget before downstream arguments may treat the native family as
adequate for the dissolving family.

The central object is the \emph{adequacy residual}.  Informally, before the
formal definition in \Cref{sec:adequacy-residual}, it is the Schur complement of
the native block in the joint native-dissolving currency:
\[
  \Xi_C(D\mid L)
  =
  K_{DD}-K_{DL}K_{LL}^{\dagger}K_{LD}.
\]
Here \(K_{LL}\) is the native block currency, \(K_{DD}\) is the dissolving block
currency, and \(K_{DL},K_{LD}\) are the cross currencies built from the two probe
families and the audit energy.  Thus \(\Xi_C(D\mid L)\) is the residual left
after the part of the dissolving response explained through the native block has
been subtracted.

The paper develops this residual through a sequence of standard
finite-dimensional operator identities, organized around the role the residual
plays in budgeted adequacy.  First, \Cref{sec:projection} identifies
\(\Xi_C(D\mid L)\) with the dissolving readout after projection away from the
energy-scaled native probe range.  This projection identity gives the Schur
complement its geometric interpretation and shows positivity directly.  Next,
\Cref{sec:optimal-exact} proves the optimal residual identity: among residual
currencies produced by explanation maps from native responses to dissolving
responses, the adequacy residual is Loewner-minimal.  The same section records
the exact-adequacy equivalences, including the zero-residual condition
\(\Xi_C(D\mid L)=0\), kernel containment on the legal quotient, and the
existence of an explanation map \(A\) with \(D_0=AL_0\).  These equivalences are
the finite ``no leftover'' tests for the framework.

The later sections record how the residual behaves under the operations needed
to use it.  \Cref{sec:promotion} shows that native and residual budgets combine
through the optimal native explanation to produce a dissolving budget; this is
the adequacy-promotion step.  \Cref{sec:chain-rule} proves the chain rule under
strict native extension: when the native family is enlarged, a conditional
Schur term is subtracted, so the residual cannot increase under the stated
hypotheses.  \Cref{sec:data-processing} gives the data-processing and transfer
identities: post-processing the dissolving probes conjugates the residual,
coarsening the native probes can only enlarge it under the comparison
hypotheses, and exact or defective bridges transport it with the corresponding
congruence or paid defect.  Finally, \Cref{sec:blind-spot} turns residual-budget
failure into a witness statement: a budget violation is equivalent to an
explicit recombination vector whose quadratic excess is positive.

The present paper is self-contained as a Schur-complement residual paper:
its hypotheses, definitions, identities, and witness theorem are stated
in the finite operator-theoretic setting just described.  An application
of the calculus to a layer-dissolving framework is given in
\citet{Tsiokos2026NeedleKiller}, which readers may consult for that use.

The scope is deliberately narrow.  The paper does not claim that every raw
substrate sees every dissolving probe, nor that the residual vanishes
generically; exact adequacy is a hypothesis to be verified, not a default
state.  It also does not transfer a residual computed from a coarse view to an
intrinsic fine view without an audited bridge: such transfer requires the exact
or defective bridge machinery developed later.  A residual bound at one
accepted finite stage is not treated as predictive adequacy across an indexed
family.  Finally, the paper does not claim novelty for the classical
constructions it uses, such as Schur complements, Moore--Penrose
pseudoinverses, projection identities, and Douglas-style factorization.  The
contribution is to package the Schur-complement residual as an audited currency
with projection, optimality, exactness, extension, data-processing, transfer,
and blind-spot witness laws.

The organization follows this progression.  \Cref{sec:setup} fixes the standing
hypotheses: the finite-dimensional carrier, the audit energy \(C\), the legal
energy quotient \(E_C\), and null-mode legality for probes.  \Cref{sec:currency}
develops probe-family currency and minimum-spend duality, while
\Cref{sec:adequacy-residual} gives the formal definition of
\(\Xi_C(D\mid L)\) and its block currencies.  The main calculus then proceeds
through projection, optimality and exactness, promotion, chain rule, data
processing and transfer, and the blind-spot witness theorem in
\Cref{sec:projection,sec:optimal-exact,sec:promotion,sec:chain-rule,sec:data-processing,sec:blind-spot}.
The discussion and conclusion close by placing the residual calculus back in
the framework that motivates it, without adding new hypotheses or applications.
 \section{Standing Hypotheses: The Legal Energy Quotient}
\label{sec:setup}

This section fixes the standing hypotheses used by every later definition and
identity in the paper.  The data are a finite-dimensional carrier, a positive
semidefinite audit energy, the legal energy quotient determined by the kernel
of that energy, and a null-mode legality condition for probe families.  Once
stated here, these hypotheses are assumed throughout the residual calculus.

\begin{definition}[Audit energy]\label{def:xi:audit-energy}
Let \(E\) be a finite-dimensional real or complex Hilbert space.  An
\emph{audit energy} on \(E\) is a self-adjoint positive-semidefinite operator
\(C=C^*\succeq 0\) on \(E\).  The audit cost of a carrier vector \(u\in E\) is
the quadratic form \(u^*Cu\), and the kernel \(\Ker C\) consists of the
zero-audit-cost directions.
\end{definition}

The form may be semidefinite, so directions in \(\Ker C\) cannot be treated as
ordinary finite-energy degrees of freedom without an additional quotient step.

\begin{definition}[Legal energy quotient]\label{def:xi:legal-energy-quotient}
The \emph{legal energy quotient} of \((E,C)\) is the orthogonal complement
\[
  E_C=(\Ker C)^\perp\subseteq E,
\]
together with the restricted audit form
\[
  C_0=C|_{E_C}>0
\]
which is strictly positive on \(E_C\).  The \emph{Moore--Penrose pseudoinverse}
\(C^\dagger\), following
\citet{Penrose1955GeneralizedInverse,BenIsraelGreville2003GeneralizedInverses},
acts as \(C_0^{-1}\) on \(E_C\) and as zero on \(\Ker C\).
\end{definition}

All finite-currency formulas in the paper are read on this legal subspace.
Equivalently, quadratic response expressions involving \(C^\dagger\) are
shorthand for first passing to \(E_C\), applying \(C_0^{-1}\), and then viewing
the result back in the original coordinates.

\begin{definition}[Null-legal probe family]\label{def:xi:null-legal-probe-family}
A linear map \(F\colon E\to Y\) is a \emph{null-legal probe family} when it
vanishes on the zero-audit-cost directions:
\[
  F|_{\Ker C}=0.
\]
For such an \(F\), its \emph{restriction to the legal subspace} is
\[
  F_0=F|_{E_C}\colon E_C\to Y.
\]
The native and dissolving probe families used throughout the paper are taken to
be null-legal: their restrictions are written \(L_0\) and \(D_0\).
\end{definition}

Null-mode legality says that the probe descends cleanly through the quotient:
no response is assigned to a direction whose audit cost is zero.

This condition is part of the adequacy record, not a disposable regularity
assumption.  If a probe sees a zero-energy direction, then the variational
capacity associated to that response is infinite, even when a formal
pseudoinverse expression might appear finite after the zero mode has been
silenced algebraically.  Throughout the paper, every probe family that enters a
currency matrix or an adequacy residual is taken to be null-legal from the
start.  The assumption is data of the probe family, checked before forming the
currency, rather than a condition repaired after the residual has been written
down.

Consequently, every later currency and residual is computed on \(E_C\) with
\(C_0\) and the restricted probes.  The convention for \(C^\dagger\) is the
canonical Moore--Penrose convention just described; the distinction between the
pseudoinverse on \(E\) and the inverse on \(E_C\) matters only as bookkeeping,
not as a separate mathematical choice in the body of the paper.  With the legal
quotient fixed, \Cref{sec:currency} develops the probe-family currency and
minimum-spend duality used throughout.  \Cref{sec:adequacy-residual} then
defines the adequacy residual \(\Xi_C(D\mid L)\) from the native, dissolving,
and cross currency blocks computed under these standing hypotheses.
 \section{Currency and Minimum Spend}
\label{sec:currency}

With the legal quotient and null-mode legality of \Cref{sec:setup} in place,
every probe family produces a currency matrix recording its response per unit
of audit cost.  The minimum-spend formula completes the duality: it prices the
least audit cost needed to realize a desired response, while the off-range
branch records that responses outside the represented range are unattainable at
finite cost.  This is the accounting language used in every later residual
identity.

\begin{definition}[Probe-family currency]\label{def:xi:probe-family-currency}
Let \(F\colon E\to Y\) be a null-legal probe family in the sense of
\Cref{def:xi:null-legal-probe-family}, with restriction \(F_0=F|_{E_C}\).
Its \emph{probe-family currency} is the positive semidefinite operator on \(Y\)
\[
  K_F = F_0C_0^{-1}F_0^* = FC^{\dagger}F^*,
\]
where the right-hand expression uses the Moore--Penrose pseudoinverse
\(C^\dagger\) of \Cref{def:xi:legal-energy-quotient}.  For a recombination
\(y\in Y\), the variational capacity of the scalar probe \(y^*F\) is
\[
  \Capacity_C(y^*F)=y^*K_Fy.
\]
\end{definition}

Thus \(K_F\) is the compliance matrix associated with \(F\): it records how much
response the probe family can produce per audit dollar in each recombined
direction of the response space.

\begin{theorem}[Minimum-spend duality]\label{thm:xi:minimum-spend-duality}
For \(z\in Y\), define
\[
  \Cost_C(z;F)
  =
  \inf\bigl\{\langle u,Cu\rangle:\ Fu=z\bigr\}.
\]
Then
\[
  \Cost_C(z;F)
  =
  \begin{cases}
    z^*K_F^\dagger z,
      & z\in \Ran F_0=\Ran K_F,\\
    +\infty,
      & z\notin \Ran F_0.
  \end{cases}
\]
\end{theorem}

\begin{proof}
Because all probes entering the currency are null-legal, the optimization may
be read on \(E_C\).  Set \(v=C_0^{1/2}u\).  Since \(C_0>0\) on \(E_C\), this
substitution is legitimate, the audit cost becomes the ordinary squared norm
\(\|v\|^2\), and the constraint \(Fu=z\) becomes
\[
  F_0C_0^{-1/2}v=z.
\]
Writing \(T=F_0C_0^{-1/2}\), the problem is therefore the standard
least-norm problem of minimizing \(\|v\|^2\) subject to \(Tv=z\).  If
\(z\in\Ran T\), the minimum is \(z^*(TT^*)^\dagger z\); if
\(z\notin\Ran T\), the admissible set is empty and the cost is \(+\infty\).
Since
\[
  TT^*=F_0C_0^{-1}F_0^*=K_F,
\]
the finite branch is \(z^*K_F^\dagger z\).  The range condition can be written
as \(z\in\Ran T=\Ran F_0=\Ran K_F\), using invertibility of \(C_0^{-1/2}\) on
\(E_C\) and the finite-dimensional identity \(\Ran(TT^*)=\Ran T\).\footnote{Mechanized in Lean as \texttt{SixBirdsNeedles.Xi.Currency.xiMinimumSpendDuality}; the Lean theorem proves the off-range \((+\infty)\) branch under the finite-range pseudoinverse identity being established here as a least-norm calculation.  The full statement is proved here.  See \Cref{subsec:paper2-coverage-table}.}
\end{proof}

The theorem gives the variational meaning of \(K_F\).  As a compliance currency,
it measures the response made available by the probe family, while its
pseudoinverse prices the least audit spend needed to realize a represented
response.  If \(z\) lies outside \(\Ran F_0\), no vector of finite audit spend
can realize it through \(F\).  The remaining sections build on this language:
\Cref{sec:adequacy-residual} constructs the joint native-dissolving currency
matrix and defines the adequacy residual as the Schur complement of the native
block.
 \section{Definition of the Adequacy Residual}
\label{sec:adequacy-residual}

This section gives the central object of the paper.  With the standing
hypotheses of \Cref{sec:setup} and the probe-family currency of
\Cref{sec:currency} in place, the joint native-dissolving currency matrix
decomposes into four blocks.  The adequacy residual is the Schur complement of
the native block
\citep{Schur1917Potenzreihen,Haynsworth1968Inertia,Albert1969PSDPartitioned,Cottle1974Manifestations,Zhang2005SchurApplications,HornJohnson2013MatrixAnalysis}:
the leftover content of the dissolving block after the native
block has explained as much as it can.

\begin{definition}[Native, dissolving, and joint block currencies]\label{def:xi:block-currencies}
Let \(L\colon E\to Y\) and \(D\colon E\to Z\) be null-legal probe families (in
the sense of \Cref{def:xi:null-legal-probe-family}), with restrictions
\(L_0,D_0\).  The \emph{native, dissolving, and cross currencies} are
\[
  K_{LL}=L_0C_0^{-1}L_0^*,
  \qquad
  K_{DL}=D_0C_0^{-1}L_0^*,
\]
\[
  K_{LD}=L_0C_0^{-1}D_0^*,
  \qquad
  K_{DD}=D_0C_0^{-1}D_0^*.
\]
Here \(K_{LL}=K_F\) for \(F=L\) and \(K_{DD}=K_F\) for \(F=D\) in the
probe-family currency of \Cref{def:xi:probe-family-currency}.  Equivalently,
\(K_{LL}\) and \(K_{DD}\) are positive semidefinite operators on \(Y\) and
\(Z\) respectively; \(K_{DL}\) and \(K_{LD}=K_{DL}^*\) are the corresponding
rectangular cross operators.  The \emph{joint currency matrix} of the combined
family \((L,D)\) is
\[
  K(D,L)
  =
  \begin{pmatrix}
    K_{LL} & K_{LD} \\
    K_{DL} & K_{DD}
  \end{pmatrix}.
\]
\end{definition}

The convention is that \(L\) is the measurement family available at the
original layer, while \(D\) is the measurement family available after passing to
the richer layer.

\begin{definition}[Adequacy residual]\label{def:xi:adequacy-residual}
The adequacy residual of \(D\) relative to \(L\) is
\[
  \boxed{\Xi_C(D\mid L)=K_{DD}-K_{DL}K_{LL}^\dagger K_{LD}.}
\]
\end{definition}

The residual is the Schur complement of the native block inside the joint
currency matrix of \Cref{def:xi:block-currencies}.  In the standard
\(2\times2\) block Schur-complement decomposition,
\(\Xi_C(D\mid L)\) is the \((2,2)\) Schur complement: \(K_{DD}\) minus the part
of the dissolving block explained through the native block by the cross term
\(K_{DL}K_{LL}^\dagger K_{LD}\).  If \(K_{LL}\) is invertible on the native
response space \(Y\), this is the ordinary inverse formula.  The Moore--Penrose
pseudoinverse is what keeps the same expression meaningful when the native
currency is rank deficient.

This is exactly the algebraic operation needed for leftover currency.  The
native block first accounts for every component of the dissolving response that
can be explained through native coordinates; the residual is what remains.
Equivalently, \(\Xi_C(D\mid L)\) is the unexplained dissolving currency under
the native block's optimal explanation.  \Cref{sec:projection} makes this
geometric by identifying the residual with a projected dissolving readout, and
\Cref{sec:optimal-exact} shows that it is Loewner-minimal among residuals
arising from explanation maps.

The joint currency matrix is positive semidefinite whenever the audit energy
is: it is the currency of the combined probe family \((L,D)\) on the legal
subspace.  The Schur complement of this positive semidefinite block matrix is
therefore positive semidefinite under the standing hypotheses.  Consequently
\(\Xi_C(D\mid L)\) is a positive operator on \(Z\), and so it is a valid
currency in the sense of \Cref{sec:currency}.  The remaining sections develop
the calculus of this positive residual.
 \section{Projection Theorem}
\label{sec:projection}

The projection theorem gives the geometric meaning of the adequacy residual.
After rescaling by the square root of the audit energy, the native and
dissolving probes become ordinary linear maps in energy-scaled coordinates.  In
those coordinates, the residual is the dissolving readout after projecting away
the part already explained by the native probe range.  Positive
semidefiniteness of \(\Xi_C(D\mid L)\) follows directly from this formulation.

\begin{definition}[Energy-scaled probes and native-explanation projection]\label{def:xi:energy-scaled-probes}
The \emph{energy-scaled native and dissolving probes} are
\[
  T_L=L_0C_0^{-1/2},
  \qquad
  T_D=D_0C_0^{-1/2}.
\]
Here \(C_0^{1/2}\) is the unique positive square root of \(C_0\) on \(E_C\)
\citep{Higham2008Functions,Bhatia2007PositiveDefinite}.  The
\emph{native-explanation projection} \(P_L\) is the orthogonal projection in
\(E_C\) onto
\[
  \Ran(T_L^*)\subseteq E_C,
\]
the native probe range in the energy-scaled coordinate system.
\end{definition}

The maps \(T_L\) and \(T_D\) record the same response data as \(L_0\) and
\(D_0\), but expressed against the energy-orthonormal coordinates determined by
\(C_0\).  The range defining \(P_L\) is not the range of \(L_0\) in the original
response coordinates.

\begin{theorem}[Projection identity]\label{thm:xi:projection-identity}
\[
  \boxed{\Xi_C(D\mid L)=T_D(I-P_L)T_D^*\succeq 0.}
\]
\end{theorem}

\begin{proof}
In terms of the energy-scaled probes, the four currency blocks are
\[
  K_{LL}=T_LT_L^*,
  \qquad
  K_{DD}=T_DT_D^*,
  \qquad
  K_{DL}=T_DT_L^*,
  \qquad
  K_{LD}=T_LT_D^*.
\]
The standard Moore--Penrose projection identity
\citep{Penrose1956BestApprox,Halmos1982HilbertSpace} says that
\[
  T_L^*(T_LT_L^*)^\dagger T_L=T_L^*K_{LL}^\dagger T_L
\]
is the orthogonal projection onto \(\Ran(T_L^*)\), namely \(P_L\).  Therefore
\[
  K_{DL}K_{LL}^\dagger K_{LD}
  =
  T_DT_L^*K_{LL}^\dagger T_LT_D^*
  =
  T_D P_L T_D^*.
\]
Substituting this into the definition of the residual gives
\[
  \Xi_C(D\mid L)
  =
  K_{DD}-T_D P_L T_D^*
  =
  T_D T_D^*-T_D P_L T_D^*
  =
  T_D(I-P_L)T_D^*.
\]
Since \(I-P_L\) is the orthogonal projection onto the complementary subspace,
it is positive semidefinite, and conjugation by \(T_D\) preserves positive
semidefiniteness.\footnote{Mechanized in Lean as \texttt{SixBirdsNeedles.Xi.Projection.projectionIdentity}; the Lean theorem proves the algebraic Schur-complement identity \(\Xi_C(D\mid L)=T_D(I-P_L)T_D^*\) under \(T_L\), \(T_D\), and \(P_L\) being exposed as free parameters under four algebraic hypotheses replacing the paper's constructive square-root and projection machinery.  The full statement is proved here.  See \Cref{subsec:paper2-coverage-table}.}
\end{proof}

The projection identity shows that \(\Xi_C(D\mid L)\) is the dissolving readout
after the part explained by native probes has been removed in the energy-dual
geometry.  Equivalently, it measures the dissolving response against the
orthogonal complement of the native probe range in \(E_C\)-scaled coordinates.
The residual is therefore a true projection residual, not merely an arbitrary
subtraction of two positive operators.

This formulation makes positivity manifest and supplies the ``removed by native
explanation'' reading used in \Cref{sec:optimal-exact}.  It also gives the
algebraic basis for the data-processing identities of
\Cref{sec:data-processing}: when dissolving probes are post-processed, the
projected readout is conjugated by the post-processing map.  The next section
quantifies optimality and exactness for the same residual.
 \section{Optimal Residual and Exact Adequacy}
\label{sec:optimal-exact}

This section answers two questions about the adequacy residual.  Among all
linear explanation maps \(A\colon Y\to Z\) that try to factor the dissolving
probe \(D\) through the native probe \(L\), which map minimizes the residual
currency?  And when is the minimum residual equal to zero?  The first answer is
the optimal residual identity, which shows that \(\Xi_C(D\mid L)\) is
Loewner-minimal across choices of \(A\).  The second is exact adequacy: four
equivalent ways of saying that no dissolving currency is left over.

\begin{definition}[Optimal native explanation map]\label{def:xi:optimal-explanation-map}
Given native and dissolving probe families \(L,D\) and their block currencies
\(K_{LL},K_{DL}\) of \Cref{def:xi:block-currencies}, the \emph{optimal native
explanation map} is
\[
  A_*=K_{DL}K_{LL}^\dagger\colon Y\to Z.
\]
When \(K_{LL}\) is invertible on \(Y\), \(A_*\) reduces to the ordinary
regression-style \(K_{DL}K_{LL}^{-1}\), the matrix-form Gauss--Markov /
generalized-least-squares explanation map
\citep{Aitken1935LeastSquares,Rao1973LinearInference}.
\end{definition}

This is the linear map that best approximates the dissolving block by
composition through the native block, in the Loewner sense made precise below.

\begin{theorem}[Optimal residual identity]
\label{thm:xi:optimal-residual-identity}
For every linear map \(A\colon Y\to Z\),
\[
  \boxed{
  (D_0-AL_0)C_0^{-1}(D_0-AL_0)^*
  =
  \Xi_C(D\mid L)
  +(A-A_*)K_{LL}(A-A_*)^*.}
\]
Consequently, \(\Xi_C(D\mid L)\) is the Loewner-minimal residual currency
among all attempts to factor \(D\) through \(L\).
\end{theorem}

\begin{proof}
Expanding the left-hand side by linearity gives
\[
  K_{DD}-AK_{LD}-K_{DL}A^*+AK_{LL}A^*.
\]
Write \(A=A_*+B\), where \(B=A-A_*\).  The Moore--Penrose identities, together
with the block-range fact that \(K_{DL}\) vanishes on \(\Ker K_{LL}\), give
\[
  A_*K_{LL}
  =
  K_{DL}K_{LL}^\dagger K_{LL}
  =
  K_{DL}
\]
and the corresponding adjoint identity for \(K_{LD}A_*^*\).  Substituting
\(A_*+B\) into the expansion, the linear terms in \(B\) telescope, while the
\(A_*\)-part leaves
\[
  K_{DD}-K_{DL}K_{LL}^\dagger K_{LD}
  =
  \Xi_C(D\mid L).
\]
The remaining quadratic term is \(BK_{LL}B^*\), which is
\((A-A_*)K_{LL}(A-A_*)^*\).  Since this correction is positive semidefinite for
every \(A\), the displayed residual is always Loewner-above
\(\Xi_C(D\mid L)\), with equality at \(A=A_*\).\footnote{Mechanized in Lean as \texttt{SixBirdsNeedles.Xi.OptimalResidual.optimalResidualIdentity}; the Lean theorem proves the Loewner-domination consequence under the paper-algebraic identity \(\mathrm{residualQuad}\,A=\mathrm{adequacyResidual}+\mathrm{penalty}(A-A_*)\) and a PSD hypothesis on \(K_{LL}\); the identity itself is established here by the expansion just given.  The full statement is proved here.  See \Cref{subsec:paper2-coverage-table}.}
\end{proof}

\begin{theorem}[Exact adequacy]\label{thm:xi:exact-adequacy}
The following are equivalent:
\begin{enumerate}[label=(\roman*)]
  \item \(\Xi_C(D\mid L)=0\);
  \item \(D_0=A_*L_0\);
  \item there exists a linear map \(A\colon Y\to Z\) such that \(D_0=AL_0\);
  \item \(\Ker L_0\subseteq \Ker D_0\).
\end{enumerate}
\end{theorem}

\begin{proof}
First assume (i).  Applying the optimal residual identity at \(A=A_*\) gives
\[
  (D_0-A_*L_0)C_0^{-1}(D_0-A_*L_0)^*
  =
  \Xi_C(D\mid L).
\]
If the residual is zero, then the left-hand side is the zero positive
semidefinite operator.  Since \(C_0^{-1}\) is strictly positive on \(E_C\), this
forces \(D_0-A_*L_0=0\), proving (ii).  The implication (ii) implies (iii) is
immediate by taking \(A=A_*\).  If (iii) holds and \(v\in\Ker L_0\), then
\(D_0v=AL_0v=0\), so \(v\in\Ker D_0\); hence (iv) follows.

It remains to see that (iv) gives zero residual.  In finite dimension, the
kernel containment \(\Ker L_0\subseteq\Ker D_0\) is exactly the factorization
criterion for the existence of a linear map \(A\colon Y\to Z\) with
\(D_0=AL_0\), in the range-inclusion / factorization tradition of
\citet{Douglas1966Majorization} and \citet{FillmoreWilliams1971OperatorRanges}.  Substituting this factorization into the optimal residual
identity gives
\[
  0
  =
  \Xi_C(D\mid L)
  +(A-A_*)K_{LL}(A-A_*)^*.
\]
Both summands are positive semidefinite, so each must vanish; in particular
\(\Xi_C(D\mid L)=0\).  Equivalently, by the projection theorem, the same
condition can be read as
\(\Xi_C(D\mid L)=0\) iff \(T_D=T_DP_L\), iff
\(\Ran(T_D^*)\subseteq\Ran(T_L^*)\), which is the energy-scaled form of
factorization through \(L_0\).\footnote{Mechanized in Lean as \texttt{SixBirdsNeedles.Xi.OptimalResidual.exactAdequacy}; the Lean theorem proves the trivial direction chains (ii)\(\to\)(iii)\(\to\)(iv) in the four-way equivalence under the hard directions (i)\(\leftrightarrow\)(ii) and (iv)\(\to\)(ii) being elevated as paper-citation hypotheses; the full cycle is established here using the Douglas--Fillmore-Williams factorization criterion.  The full statement is proved here.  See \Cref{subsec:paper2-coverage-table}.}
\end{proof}

The exact-adequacy theorem is the paper's finite ``no leftover'' test.  It says
that zero residual, exact factorization through the native family, existence of
some explanation map, and kernel containment are the same condition on the
legal quotient.  Together with the optimal residual identity, it makes
\(\Xi_C(D\mid L)\) the canonical leftover currency selected by the
Schur-complement framework, not merely one possible residual.  The next
sections develop its behavior under operations: \Cref{sec:promotion} combines
native and residual budgets into a dissolving budget, while
\Cref{sec:chain-rule,sec:data-processing} track how the residual changes under
probe-family extension, coarsening, post-processing, and transfer.
 \section{Promotion from Native Membrane to Dissolving Membrane}
\label{sec:promotion}

The previous sections built the adequacy residual, gave its projection
interpretation, and characterized when it vanishes.  The present section records
the operational consequence.  When a native budget and a residual budget are
both available, they combine through the optimal native explanation to produce
a dissolving budget.  This is the promotion step from a native membrane budget
record to a dissolving one, and it is the form used at the layer-dissolving
endpoint of \citet{Tsiokos2026NeedleKiller}.

Suppose first that the native currency is controlled by a positive operator
\(\Theta_Y\) on \(Y\):
\[
  K_{LL}\preceqq \Theta_Y.
\]
This is the \emph{native budget}.  Suppose also that the adequacy residual is
controlled by a positive operator \(\Xi_Z\) on \(Z\):
\[
  \Xi_C(D\mid L)\preceqq \Xi_Z.
\]
This is the \emph{residual budget}.  The two hypotheses are independent: the
first prices what the native probes already explain, while the second prices
what remains after native explanation.

\begin{theorem}[Adequacy promotion]\label{thm:xi:adequacy-promotion}
Under the two budget hypotheses above,
\[
  \boxed{K_{DD}\preceqq A_*\Theta_YA_*^*+\Xi_Z.}
\]
In particular, if \(\Xi_C(D\mid L)=0\), the native budget transfers exactly to
the dissolving family through \(A_*\).
\end{theorem}

\begin{proof}
The Schur identity obtained from the optimal residual identity at
\(A=A_*\) gives
\[
  K_{DD}=A_*K_{LL}A_*^*+\Xi_C(D\mid L).
\]
Apply the native budget and conjugate by \(A_*\); Loewner order is monotone
under conjugation, so
\[
  A_*K_{LL}A_*^*\preceqq A_*\Theta_YA_*^*.
\]
The residual budget gives \(\Xi_C(D\mid L)\preceqq\Xi_Z\).  Adding the two
Loewner inequalities yields the displayed promotion bound.  If exact adequacy
holds, then \(\Xi_C(D\mid L)=0\), so no residual budget is needed and the
native budget transfers through \(A_*\) alone.\footnote{Mechanized in Lean as \texttt{SixBirdsNeedles.Xi.Promotion.adequacyPromotion}; the Lean theorem proves the Loewner promotion bound \(K_{DD}\preceqq A_*\Theta_Y A_*^*+\Xi_Z\) under the adequacy-interface algebraic identity and the public Loewner-monotonicity infrastructure; the identity is the elevated hypothesis.  The full statement is proved here.  See \Cref{subsec:paper2-coverage-table}.}
\end{proof}

Operationally, adequacy promotion is the bridge result invoked at the
layer-dissolving endpoint of \citet{Tsiokos2026NeedleKiller}.  The optimal explanation map \(A_*\)
plays the role of a best linear predictor of the dissolving response from the
native response, and adequacy promotion is the corresponding budget transfer;
the analogous orthogonality-principle and best-linear-prediction picture is
developed in \citet{KailathSayedHassibi2000LinearEstimation}.
There, a native budget and a residual budget are converted into a dissolving
budget through the Xi interface, and the resulting dissolving currency can then be routed into a
common target space by the maps belonging to that application.  From the
standpoint of this paper, promotion is the simplest direct use of the optimal
residual identity: Schur decomposition plus two Loewner bounds gives the
dissolving budget.  For the endpoint application, see
\citet{Tsiokos2026NeedleKiller}.
 \section{Strict Native Extension}
\label{sec:chain-rule}

This section asks how the adequacy residual changes when the native probe
family \(L\) is enlarged by adding new probes.  The answer is a Schur-complement
chain rule: extending the native family subtracts a conditional residual term
built from the added probes.  Consequently the residual can only shrink under
strict native extension.  If the added probes are already linear functions of
the original native family, the residual is unchanged.

\begin{definition}[Extended native family and conditional currencies]\label{def:xi:extended-native-family}
Let \(M\colon E\to W\) be an additional null-legal native probe family.  The
\emph{extended native family} is
\[
  L^+
  =
  \begin{bmatrix}
    L\\
    M
  \end{bmatrix}
  \colon E\to Y\oplus W,
  \qquad
  L^+u=(Lu,Mu),
\]
which is null-legal because \(L\) and \(M\) both vanish on \(\Ker C\).  The
\emph{conditional native currency of \(M\) given \(L\)} is the Schur complement
\[
  K_{MM\mid L}
  =
  K_{MM}-K_{ML}K_{LL}^\dagger K_{LM},
\]
also known as the shorted operator of \(K_{MM}\) by the \(L\)-block
\citep{AndersonWN1971Shorted,AndersonTrapp1975ShortedII}.  The
\emph{conditional dissolving-extension cross currency} is
\[
  K_{DM\mid L}
  =
  K_{DM}-K_{DL}K_{LL}^\dagger K_{LM},
  \qquad
  K_{MD\mid L}=K_{DM\mid L}^*.
\]
The conditional joint currency block is positive semidefinite, so in
particular \(K_{MM\mid L}\succeq 0\) on \(W\); the rectangular cross
block \(K_{DM\mid L}\) is paired with its adjoint \(K_{MD\mid L}\).
\(K_{MM\mid L}\) vanishes precisely when the added probes \(M\) are
already explained by \(L\) on the legal quotient.
\end{definition}

\begin{theorem}[Chain rule for adequacy residuals]\label{thm:xi:chain-rule}
\[
  \boxed{
  \Xi_C(D\mid L^+)
  =
  \Xi_C(D\mid L)
  -
  K_{DM\mid L}K_{MM\mid L}^\dagger K_{MD\mid L}.}
\]
In particular,
\[
  \Xi_C(D\mid L^+)\preceqq \Xi_C(D\mid L).
\]
\end{theorem}

\begin{proof}
Consider the joint currency matrix with rows and columns indexed by \(L\),
\(M\), and \(D\).  Its diagonal blocks are \(K_{LL}\), \(K_{MM}\), and
\(K_{DD}\), with the corresponding cross-currency blocks off the diagonal.
Taking the Schur complement of the \(L\)-block gives the conditional joint
currency of \((M,D)\) given \(L\):
\[
  \begin{pmatrix}
    K_{MM\mid L} & K_{MD\mid L}\\
    K_{DM\mid L} & K_{DD\mid L}
  \end{pmatrix},
  \qquad
  K_{DD\mid L}
  =
  K_{DD}-K_{DL}K_{LL}^\dagger K_{LD}
  =
  \Xi_C(D\mid L).
\]
Computing \(\Xi_C(D\mid L^+)\) means taking the Schur complement of the full
\((L,M)\)-native block.  By the iterated Schur-complement identity, also known
as the Crabtree--Haynsworth quotient formula
\citep{CrabtreeHaynsworth1969Quotient,Ando2005SchurOperator}, this is the same
as first
conditioning on \(L\) and then taking the Schur complement of the remaining
\(M\)-block inside the conditional \((M,D)\) matrix.  Therefore
\[
  \Xi_C(D\mid L^+)
  =
  \Xi_C(D\mid L)
  -
  K_{DM\mid L}K_{MM\mid L}^\dagger K_{MD\mid L}.
\]
The subtracted term is positive semidefinite, because it is the usual
cross-block--pseudoinverse--cross-block term associated to a positive
conditional Schur block.  Hence
\(\Xi_C(D\mid L^+)\preceqq \Xi_C(D\mid L)\).\footnote{Mechanized in Lean as \texttt{SixBirdsNeedles.Xi.StrictExtension.chainRule}; the Lean theorem proves the Loewner-domination consequence under the elevated Schur chain-rule identity and PSD-of-conditional-Schur hypotheses; the identity itself is derived here from the Crabtree--Haynsworth iterated Schur formula.  The full statement is proved here.  See \Cref{subsec:paper2-coverage-table}.}
\end{proof}

\begin{corollary}[Same-family saturation]\label{cor:xi:same-family-saturation}
If \(M=BL\) for some linear map \(B\colon Y\to W\), then
\[
  \Xi_C(D\mid [L;M])=\Xi_C(D\mid L).
\]
\end{corollary}

\begin{proof}
If \(M=BL\), then \(M_0=BL_0\) on the legal quotient.  Thus
\[
  K_{ML}=BK_{LL},
  \qquad
  K_{LM}=K_{LL}B^*,
  \qquad
  K_{MM}=BK_{LL}B^*.
\]
Substitution into the conditional \(M\)-currency gives
\[
  K_{MM\mid L}
  =
  BK_{LL}B^*-BK_{LL}K_{LL}^\dagger K_{LL}B^*
  =
  0,
\]
using \(K_{LL}K_{LL}^\dagger K_{LL}=K_{LL}\).  Similarly,
\[
  K_{DM}=K_{DL}B^*,
  \qquad
  K_{DM\mid L}
  =
  K_{DL}B^*-K_{DL}K_{LL}^\dagger K_{LL}B^*
  =
  0.
\]
The chain-rule subtraction therefore vanishes, so
\(\Xi_C(D\mid L^+)=\Xi_C(D\mid L)\), equivalently
\(\Xi_C(D\mid [L;M])=\Xi_C(D\mid L)\).
\end{proof}

The chain rule is the residual analogue of partial regression.  Adding new
native probes can only subtract from the residual budget, and the subtraction
depends only on the component of the new probes that remains after conditioning
on \(L\).  Same-family saturation makes the boundary precise: adding probes
that are already linear combinations of \(L\) does not change the residual.
\Cref{sec:data-processing} turns next to the dual operations of post-processing
the dissolving probes and coarsening the native probes.
 \section{Data Processing and Transfer}
\label{sec:data-processing}

This section records how the adequacy residual behaves under the natural
operations on probe families and carriers.  On the dissolving side, a response
map simply post-processes the residual by congruence.  On the native side,
coarsening the available probes can only enlarge the residual, because the
explanation subspace has become smaller.  Finally, bridges between legal
energy quotients transport residuals exactly when the bridge is exact, and with
an explicit defect term when the dissolving part of the bridge has an error.
Together these four laws turn \(\Xi_C(D\mid L)\) from a single Schur complement
into a calculus.

\begin{theorem}[Dissolving post-processing]
\label{thm:xi:dissolving-post-processing}
For a response map \(F\colon Z\to Z'\),
\[
  \boxed{\Xi_C(FD\mid L)=F\Xi_C(D\mid L)F^*.}
\]
\end{theorem}

\begin{proof}
Replace \(D\) by \(FD\) in the block currencies of
\Cref{def:xi:adequacy-residual}.  The native block is unchanged, while every
dissolving block is transformed by \(F\):
\[
  K_{(FD)(FD)}=FK_{DD}F^*,
  \qquad
  K_{(FD)L}=FK_{DL},
  \qquad
  K_{L(FD)}=K_{LD}F^*.
\]
Substitution gives\footnote{Verified in Lean as \texttt{SixBirdsNeedles.Xi.DataProcessing.dissolvingPostProcessing}; see \Cref{subsec:paper2-coverage-table}.}
\[
  \Xi_C(FD\mid L)
  =
  FK_{DD}F^*-FK_{DL}K_{LL}^\dagger K_{LD}F^*
  =
  F\bigl(K_{DD}-K_{DL}K_{LL}^\dagger K_{LD}\bigr)F^*
  =
  F\Xi_C(D\mid L)F^*.
\]
\end{proof}

\begin{definition}[Faithful reparameterization on the native range]\label{def:xi:faithful-reparameterization}
A linear map \(A\colon Y\to Y'\) is a \emph{faithful
reparameterization on \(\Ran L_0\)} when its restriction to
the range of the native probe family is injective:
\[
  \Ker A\cap \Ran L_0=\{0\}.
\]
Equivalently, in the energy-scaled coordinates of
\Cref{def:xi:energy-scaled-probes},
\(\Ran(T_{AL}^*)=\Ran(T_L^*)\) and the native-explanation
projection is unchanged: \(P_{AL}=P_L\).  Under this
condition the Moore--Penrose comparison
\(A^*(AK_{LL}A^*)^{\dagger}A=K_{LL}^{\dagger}\) holds on the
range relevant to the Schur correction.
\end{definition}

\begin{theorem}[Native coarsening]\label{thm:xi:native-coarsening}
For a native response map \(A\colon Y\to Y'\),
\[
  \boxed{\Xi_C(D\mid AL)\succeq \Xi_C(D\mid L).}
\]
If \(A\) is a faithful reparameterization on \(\Ran L_0\)
in the sense of \Cref{def:xi:faithful-reparameterization},
equality holds.
\end{theorem}

\begin{proof}
For the energy-scaled probes of \Cref{sec:projection}, the coarsened native
probe is
\[
  T_{AL}=A T_L.
\]
Hence the explanation range becomes
\[
  \Ran(T_{AL}^*)=\Ran(T_L^*A^*)\subseteq \Ran(T_L^*).
\]
Let \(P_{AL}\) and \(P_L\) be the orthogonal projections onto these two
subspaces.  Since \(\Ran(T_{AL}^*)\subseteq\Ran(T_L^*)\), one has
\(P_{AL}\preceqq P_L\), and therefore \(I-P_{AL}\succeq I-P_L\).  The
projection identity gives
\[
  \Xi_C(D\mid AL)
  =
  T_D(I-P_{AL})T_D^*
  \succeq
  T_D(I-P_L)T_D^*
  =
  \Xi_C(D\mid L).
\]
If \(A\) is a faithful reparameterization on \(\Ran L_0\) in the sense of
\Cref{def:xi:faithful-reparameterization}, then coarsening is only a
change of coordinates on that range: \(\Ran(T_{AL}^*)=\Ran(T_L^*)\).  Thus
\(P_{AL}=P_L\), and equality follows.\footnote{Mechanized in Lean as \texttt{SixBirdsNeedles.Xi.DataProcessing.nativeCoarsening}; the Lean theorem proves the Loewner inequality \(\Xi_C(D\mid AL)\succeq \Xi_C(D\mid L)\) under the projection-comparison algebraic hypothesis on \(A^*(AK_{LL}A^*)^\dagger A\) versus \(K_{LL}^\dagger\); the reduction to comparing native-explanation subspaces is the argument here.  The full statement is proved here.  See \Cref{subsec:paper2-coverage-table}.}
\end{proof}

\begin{theorem}[Exact reducing bridge]\label{thm:xi:exact-reducing-bridge}
Suppose \(U\colon E'_{C'}\to E_C\) is an isometric energy bridge with
\[
  C'_0=U^*C_0U,
  \qquad
  L'_0=AL_0U,
  \qquad
  D'_0=BD_0U,
\]
where \(A\) is a faithful reparameterization on \(\Ran L_0\) in the
sense of \Cref{def:xi:faithful-reparameterization}.  Assume the
bridge is reducing on the probe-image subspace, so that the
transported-currency identities
\[
  K'_{L'L'}=AK_{LL}A^*,
  \qquad
  K'_{D'L'}=BK_{DL}A^*,
  \qquad
  K'_{D'D'}=BK_{DD}B^*
\]
hold.  Then
\[
  \boxed{\Xi_{C'}(D'\mid L')=B\Xi_C(D\mid L)B^*.}
\]
\end{theorem}

\begin{proof}
Substituting the transported-currency identities into the Schur-complement
definition of \(\Xi_{C'}(D'\mid L')\) gives
\[
  \Xi_{C'}(D'\mid L')
  =
  BK_{DD}B^*
  -
  BK_{DL}A^*(AK_{LL}A^*)^\dagger AK_{LD}B^*.
\]
By \Cref{def:xi:faithful-reparameterization}, faithfulness of
\(A\) on the represented native range supplies the
pseudoinverse comparison
\[
  A^*(AK_{LL}A^*)^\dagger A=K_{LL}^\dagger
\]
on the range relevant to the Schur correction.  Therefore the last display is
\[
  B\bigl(K_{DD}-K_{DL}K_{LL}^\dagger K_{LD}\bigr)B^*
  =
  B\Xi_C(D\mid L)B^*.
\]
If \(A\) is not a faithful reparameterization in the sense of
\Cref{def:xi:faithful-reparameterization}, this comparison can fail,
and the situation reverts to the coarsening direction of
\Cref{thm:xi:native-coarsening}.\footnote{Mechanized in Lean as \texttt{SixBirdsNeedles.Xi.DataProcessing.exactReducingBridge}; the Lean theorem proves the bridge residual congruence under separate transported \(K_{DD}\) and Schur-correction identities under the bridge equations; the identities are derived here by direct calculation.  The full statement is proved here.  See \Cref{subsec:paper2-coverage-table}.}
\end{proof}

\begin{theorem}[Defective dissolving bridge]\label{thm:xi:defective-bridge}
Decompose the target dissolving probe as
\[
  D'_0 = \widetilde D'_0 + R,
  \qquad
  \widetilde D'_0 = BD_0U,
\]
and suppose the native bridge is exact and that the exact source
probe \(\widetilde D'_0\) satisfies the hypotheses of
\Cref{thm:xi:exact-reducing-bridge}: the bridge is reducing on
the probe-image subspace, so that the transported-currency
identities \(K'_{L'L'}=AK_{LL}A^*\),
\(K'_{\widetilde D' L'}=BK_{DL}A^*\), and
\(K'_{\widetilde D' \widetilde D'}=BK_{DD}B^*\) hold.  Then for
every \(t>0\),
\[
  \boxed{
  \Xi_{C'}(D'\mid L')
  \preceqq
  (1+t)B\Xi_C(D\mid L)B^*
  +(1+t^{-1})\Xi_{C'}(R\mid L').}
\]
\end{theorem}

\begin{proof}
Apply the projection theorem on the primed side:
\[
  \Xi_{C'}(D'\mid L')=T_{D'}(I-P_{L'})T_{D'}^*.
\]
The decomposition \(D'_0 = \widetilde D'_0 + R\) splits the
energy-scaled dissolving probe as
\[
  T_{D'}=X+Y,
\]
where \(X=T_{\widetilde D'}\) is the transported contribution
from \(\widetilde D'_0=BD_0U\), and \(Y=T_R\) is the contribution
of the defect \(R\).  After projecting away the target native
explanation space, the standard Young splitting for operators
\citep{BhatiaDavis1995CauchySchwarz,Ando1995MatrixYoung} gives, for every
\(t>0\),
\[
  (X+Y)(I-P_{L'})(X+Y)^*
  \preceqq
  (1+t)X(I-P_{L'})X^*
  +(1+t^{-1})Y(I-P_{L'})Y^*.
\]
The first term is identified by
\Cref{thm:xi:exact-reducing-bridge} as \(B\Xi_C(D\mid L)B^*\),
since \(\widetilde D'_0=BD_0U\) satisfies the exact-reducing-
bridge hypotheses by the assumption above.  The second term is
precisely \(\Xi_{C'}(R\mid L')\), the residual of the defect
alone against the target native family.  This gives the
displayed bound.\footnote{Mechanized in Lean as \texttt{SixBirdsNeedles.Xi.DataProcessing.defectiveBridge}; the Lean theorem proves the defective bridge bound for every \(t>0\) under a residual decomposition and the public matrix AM--GM Loewner lemma; PSD of the transported source residual and of the defect residual are elevated hypotheses, and the Young-splitting reduction is the argument here.  The full statement is proved here.  See \Cref{subsec:paper2-coverage-table}.}
\end{proof}

These four theorems are the data-processing and transfer calculus for the
adequacy residual.  Post-processing on the dissolving side acts by congruence;
coarsening on the native side enlarges the residual; exact bridges transport
the residual by congruence on the dissolving response space; and defective
bridges add an explicit defect residual that must be paid in the target budget.
\Cref{sec:blind-spot} turns from transport to obstruction and repair: when a
residual budget fails, the failure is witnessed by an explicit recombination
vector.
 \section{Obstruction and Repair Calculus}
\label{sec:blind-spot}

The calculus so far has built the adequacy residual, identified its projection
geometry, and described how it behaves under extension, processing, and
transfer.  This section turns the same residual into a budgeted obstruction.
When a declared residual budget fails, the failure is not merely formal: it is
detected by an explicit recombination vector in the dissolving response space.
Such a vector is the blind-spot witness, and the section closes by listing the
lawful repairs available after the witness has exposed the failure.

\begin{definition}[Declared residual budget and adequacy defect]\label{def:xi:adequacy-defect}
A \emph{declared residual budget} on \(Z\) is a positive semidefinite operator
\(\Omega\succeq 0\).  The corresponding \emph{adequacy defect} of \(D\)
relative to \(L\) and \(\Omega\) is the Hermitian operator
\[
  \Delta_\Xi=\Xi_C(D\mid L)-\Omega,
\]
where \(\Xi_C(D\mid L)\) is the adequacy residual of
\Cref{def:xi:adequacy-residual}.  The budget holds when
\(\Delta_\Xi\preceqq 0\); a positive quadratic-form direction of
\(\Delta_\Xi\) is a certificate that it does not.
\end{definition}

The defect need not be positive semidefinite; only one positive quadratic-form
direction is needed to witness failure.

\begin{theorem}[Blind-spot witness]\label{thm:xi:blind-spot-witness}
\[
  \Xi_C(D\mid L)\npreceq \Omega
\]
if and only if there exists \(z\in Z\) such that
\[
  \boxed{z^*\bigl(\Xi_C(D\mid L)-\Omega\bigr)z>0.}
\]
Moreover, the maximum additive violation is \(\lambda_{\max}(\Delta_\Xi)\),
the largest eigenvalue of the adequacy defect.
\end{theorem}

\begin{proof}
The Loewner inequality \(A\preceqq B\) is equivalent to
\[
  z^*(B-A)z\geq 0
\]
for every vector \(z\).  Its negation is therefore the existence of some vector
whose quadratic form makes \(B-A\) negative, equivalently whose quadratic form
makes \(A-B\) positive.  Applying this with
\[
  A=\Xi_C(D\mid L),
  \qquad
  B=\Omega,
\]
shows that \(\Xi_C(D\mid L)\npreceq\Omega\) is exactly the existence of
\(z\in Z\) with \(z^*\Delta_\Xi z>0\).  The largest normalized value of this
quadratic form is
\[
  \sup_{z\neq 0}\frac{z^*\Delta_\Xi z}{z^*z}
  =
  \lambda_{\max}(\Delta_\Xi),
\]
by the Rayleigh-quotient principle for Hermitian operators
\citep{Bhatia1997MatrixAnalysis,HornJohnson2013MatrixAnalysis}.  Thus
\(\lambda_{\max}(\Delta_\Xi)\) is the largest single-vector additive violation.\footnote{Mechanized in Lean as \texttt{SixBirdsNeedles.Xi.Obstruction.blindSpotWitness}; the Lean theorem proves the witness equivalence \(\Xi_C(D\mid L)\npreceq\Omega\) iff there exists \(z\) with \(z^*\Delta_\Xi z>0\) under the identification of the maximum additive violation with \(\lambda_{\max}(\Delta_\Xi)\) being the Rayleigh-quotient consequence proved here.  The full statement is proved here.  See \Cref{subsec:paper2-coverage-table}.}
\end{proof}

Once a blind-spot witness has been found, the repair options are constrained.
Each lawful repair either reduces the residual by adding genuine explanatory
data, improves the bridge through which the residual is transported, or narrows
the claim being made.
\begin{enumerate}[label=(\alph*)]
  \item \emph{Strict native extension} adds new native probes and subtracts a
  positive Schur term from the residual, as in
  \Cref{thm:xi:chain-rule}.

  \item \emph{Bridge repair} improves the audited bridge data so that the
  transferred residual is lowered or the explicit defect term is reduced, as
  in \Cref{thm:xi:exact-reducing-bridge,thm:xi:defective-bridge}.

  \item \emph{Null-mode repair} annihilates or quotients away illegal
  zero-energy readouts so that the standing null-mode legality hypothesis is
  restored on a corrected carrier.

  \item \emph{Budget repair} weakens the claim by increasing \(\Omega\).  This
  may make the inequality true, but it does not reduce the blind spot itself.

  \item \emph{Nonclaim or gating} narrows the scope of the claim by excluding
  the offending response direction instead of proving adequacy for it.
\end{enumerate}
Same-family repetition saturates this list: repeating an old probe, or
repackaging the same native data, does not lower the witness because it does
not change \(\Xi_C(D\mid L)\), as shown by
\Cref{cor:xi:same-family-saturation}.  Thus a repair that merely renames the
existing native family is not a residual repair.  This catalog is the dual of
the all-six repair language of \citet{Tsiokos2026NeedleKiller}: here the same
distinction is expressed entirely in residual-budget terms.

The blind-spot witness theorem closes the calculus.  Every quantitative claim
that a residual budget fails reduces to an explicit recombination vector with
positive excess, and the excess is measured by the adequacy defect.  Combined
with the projection, optimality, exactness, chain-rule, and data-processing
identities of the preceding sections, this makes the adequacy residual a
calculable currency rather than an opaque error term.  This witness equivalence
is used inside the all-six taxonomy of \citet{Tsiokos2026NeedleKiller} to
identify failure modes precisely.
 

\section{Discussion}
\label{sec:discussion}

The paper has developed a residual calculus from one starting point: the Schur
complement
\[
  \Xi_C(D\mid L)=K_{DD}-K_{DL}K_{LL}^\dagger K_{LD}
\]
formed on the legal energy quotient.  From that definition, the preceding
sections proved the projection identity of \Cref{sec:projection}, the
optimal-residual identity and Loewner minimality across explanation maps, the
four-form exact-adequacy theorem, the adequacy-promotion step of
\Cref{thm:xi:adequacy-promotion}, the chain rule under strict native
extension, the four data-processing and transfer laws including
\Cref{thm:xi:dissolving-post-processing} and
\Cref{thm:xi:defective-bridge}, and the blind-spot witness characterization of
\Cref{sec:blind-spot}.  Together these results give a self-contained
Schur-complement-residual calculus under the legal-quotient hypotheses of
\Cref{sec:setup}.

The local mathematical tools are classical.  The construction uses Schur
complements, Moore--Penrose pseudoinverses, the projection identity
\[
  T_L^*(T_LT_L^*)^\dagger T_L
\]
as the orthogonal projection onto \(\Ran(T_L^*)\), and Douglas-style
factorization ideas behind the range-inclusion arguments for faithful
reparameterizations and bridges.  The paper does not claim novelty for those
operator-theoretic ingredients.  Its contribution is the framing: the Schur
complement is treated as an audited residual currency, with projection,
optimality, exactness, extension, data-processing, transfer, and witness laws
all stated against one energy quotient and one null-mode legality discipline.
This is why the object is not merely a block-matrix calculation but a residual
accounting device.

The phrase ``blind-spot currency'' is meant in that accounting sense.  A
nonzero residual alone says that some dissolving response is not explained by
the native family; it does not say how the failure is budgeted, transported, or
repaired.  Treating the residual as a currency makes it a first-class positive
operator with a declared budget, a witness theorem, and lawful repair
operations.  From the standpoint of this paper, the witness theorem of
\Cref{sec:blind-spot} is what gives quantitative residual claims their
operational meaning: a failed budget is equivalent to a concrete recombination
vector with positive excess.

The pseudoinverse convention is also part of the standing setup, not an
auxiliary choice made separately in each theorem.  The paper uses the
Moore--Penrose pseudoinverse on the legal quotient: \(C^\dagger\) acts as
\(C_0^{-1}\) on \(E_C\) and as zero on \(\Ker C\), and the block pseudoinverses
are the corresponding Moore--Penrose pseudoinverses of positive operators on
their response spaces.  The null-mode legality discipline of \Cref{sec:setup}
is what makes this convention meaningful for all probe families entering a
currency or residual.

An application of this calculus in a layer-dissolving framework is given in
\citet{Tsiokos2026NeedleKiller}.  In that application, the adequacy-promotion
theorem is the gate converting native and residual bounds into a dissolving
budget; the data-processing and bridge identities transport the residual
through audited bridges; and the blind-spot witness theorem supplies the
failure-detection statement inside the all-six taxonomy developed there.  The
present paper does not develop those applications.  Its task is to supply the
self-contained Schur-complement calculus that such applications can cite.

Several scope fences are therefore part of the claim.  The paper does not
assert that every raw substrate sees every dissolving probe, and it does not
assert that \(\Xi_C(D\mid L)=0\) generically; exact adequacy is a hypothesis to
verify, not a default state.  It does not promote a coarse residual view to an
intrinsic fine-grained residual without an audited bridge, and it does not
treat a finite-stage adequacy statement as predictive adequacy across an
indexed family.  Finally, it does not claim novelty for Schur complements,
pseudoinverses, projection identities, or Douglas-style factorization.  The
claim is the residual-calculus packaging under the stated hypotheses.

Several directions remain natural.  The defective dissolving bridge of
\Cref{thm:xi:defective-bridge} should admit sharper source-defect tradeoffs
under stronger structural hypotheses on the bridge map.  Predictive adequacy
ladders, deferred to appendix material in this preparation arc, extend the
chain-rule viewpoint across indexed families of native extensions and are the
next development for layered or staged settings.  The witness theorem may also
support finer repair classifications under additional record-keeping
discipline; the all-six taxonomy of \citet{Tsiokos2026NeedleKiller} is one
such refinement.
 

\section{Conclusion}
\label{sec:conclusion}

The adequacy residual \(\Xi_C(D\mid L)\) is the Schur-complement leftover
currency recording what the dissolving probe family \(D\) measures beyond what
the native probe family \(L\) explains, relative to the audit energy on the
legal quotient.  This paper developed that object as a calculus: projection,
optimal residuality, exact adequacy, promotion, strict-extension chain rule,
data processing, transfer, and blind-spot witnessing.  The result is a
self-contained Schur-complement-residual calculus on the legal energy quotient.

The calculus stands on its own as finite-dimensional operator theory under the
standing hypotheses of \Cref{sec:setup}.  An application to a layer-dissolving
framework is given in \citet{Tsiokos2026NeedleKiller}, where adequacy
promotion converts native and residual budgets into a transferred dissolving
budget.  From the standpoint of this paper, that is one use case; readers who
want only the residual calculus need not consult that work.

The residual is instance-specific.  The paper does not classify all possible
probe-family pairs \((L,D)\), does not claim exact adequacy generically, does
not claim novelty for the classical operator-theoretic constructions it uses,
and does not depend on any particular application instance.  Its theorem-level
claim is the Schur-complement-residual calculus under the stated hypotheses.

Natural next steps are correspondingly precise.  Predictive adequacy ladders
extend the chain rule across indexed native extensions; sharper defect-bridge
tradeoffs may follow from stronger involution or structural hypotheses; and
the witness theorem can support refined repair classifications under additional
record-keeping discipline.  One natural future use is to deploy the witness
theorem inside an audited framework with the all-six discipline described in
\citet{Tsiokos2026NeedleKiller}.
 
\appendix
\section{Lean Formalization Disclosure}
\label{app:formalization}

\subsection{Trust base and methodology}
\label{subsec:paper2-trust-methodology}

The Lean artifacts cited in this appendix were checked with Lean 4.28.0
\citep{deMouraUllrich2021Lean4,MathlibCommunity2020Lean}.  For
every cited theorem entry, the audited axiom closure is contained in the
standard trust base: \texttt{propext}, \texttt{Classical.choice}, and
\texttt{Quot.sound}.  The project discipline forbids \texttt{sorry},
project-local \texttt{axiom}, and \texttt{opaque} declarations in the needles
Lean corpus; \texttt{lean/manifests/trust\_base.txt} is the locked allowlist
used by the manifest checker when auditing theorem dependency closures.
Coverage is recorded per row as \emph{theorem}, \emph{partial},
\emph{definition}, or \emph{not mechanized}; these labels govern the wording
below.

\subsection{Mechanization disclosure tiers}
\label{subsec:paper2-mechanization-disclosure-tiers}

Body footnotes use a three-tier disclosure discipline keyed to the alignment
column of the per-label coverage table.

\begin{table}[H]
  \centering
  \small
  \begin{tabularx}{\linewidth}{>{\raggedright\arraybackslash}p{0.25\linewidth} >{\raggedright\arraybackslash}p{0.11\linewidth} >{\raggedright\arraybackslash}X}
    \toprule
    Alignment & Tier & Body footnote wording template \\
    \midrule
    \texttt{faithful} & T1 & Verified in Lean as \(\langle\)declaration\(\rangle\); see \Cref{subsec:paper2-coverage-table}. \\
    \texttt{narrowed-surrogate / partial} & T2 & Mechanized in Lean as \(\langle\)declaration\(\rangle\); the Lean theorem proves \(\langle\)sub-claim\(\rangle\) under \(\langle\)elevation\(\rangle\).  The full statement is proved here.  See \Cref{subsec:paper2-coverage-table}. \\
    \texttt{projection-packaged} & T3 & Tracked in the Lean formalization as \(\langle\)declaration\(\rangle\): a conditional schema over \(\langle\)record fields\(\rangle\) rather than an independent derivation, which is proved here.  See \Cref{subsec:paper2-coverage-table}. \\
    \bottomrule
  \end{tabularx}
  \caption{Paper 2 body footnote disclosure tiers.}
  \label{tab:paper2-mechanization-disclosure-tiers}
\end{table}

This tier discipline is the canonical wording rule for body footnotes: a
reader can cross-check each theorem footnote against the alignment column of
the coverage table below.

\subsection{Standing narrowings and elevated hypotheses}
\label{subsec:paper2-standing-narrowings}

Several Lean fragments use the concrete diagonal pseudoinverse
\texttt{diagonalPseudoInverse C} as the finite-matrix representative of the
paper's Moore--Penrose inverse \(C^\dagger\).  This is paper-faithful on the
legal quotient under the corresponding strict-positivity support hypotheses,
but the full quotient-level pseudoinverse theory is not mechanized in this
codebase.  When the body invokes \(C^\dagger\), the per-label disclosure below
records whether Lean proves the stated clause directly or proves a finite
surrogate under stricter positivity.

Many partial-status theorems also carry elevated algebraic hypotheses, such as
transposition symmetry, Moore--Penrose-on-support identities, or explicit
Schur-update assumptions.  The elevation is documented per row in the coverage
table; the paper supplies the ordinary mathematical proof or explanation for
the remaining clause.  These elevations should be stated in mathematical
language when relevant and should not be hidden behind manifest terminology.

The projection-packaged status flags statements whose Lean form retrieves data
from a proof-carrying record \citep{Necula1997ProofCarryingCode} or
conditional schema rather than deriving the paper theorem independently.  Such rows are traceability coverage, not
substantive Lean proof coverage.  The body may rely on the mathematical
statement, but the Lean mention must use the tier wording in
\Cref{tab:paper2-mechanization-disclosure-tiers}.

\subsection{Per-label coverage table}
\label{subsec:paper2-per-label-coverage}
\label{subsec:paper2-coverage-table}

Table~\ref{tab:paper2-lean-coverage} records the per-row Lean
coverage, alignment, and short disclosure note for each
statements-of-record row.
Table~\ref{tab:paper2-nonmechanized-coverage} lists the rows
retained with no Lean entry.

\begin{table}[H]
  \centering
  \footnotesize
  \setlength{\tabcolsep}{2.5pt}
  \renewcommand{\arraystretch}{0.94}
  \begin{tabularx}{\linewidth}{>{\raggedright\arraybackslash}p{0.37\linewidth} >{\raggedright\arraybackslash}p{0.12\linewidth} >{\raggedright\arraybackslash}p{0.15\linewidth} >{\raggedright\arraybackslash}X}
    \toprule
    Paper label & Lean coverage & Alignment & Disclosure note \\
    \midrule
    \path{thm:xi:minimum-spend-duality} & partial & needs-strengthening & Lean proves the off-range infinite-cost branch; the finite branch is paper-side. \\
    \path{def:xi:adequacy-residual} & definition & narrowed-surrogate & The Lean realization exposes the pseudoinverse as a parameter; the paper states the Moore--Penrose object. \\
    \path{thm:xi:projection-identity} & partial & narrowed-surrogate & Lean proves the algebraic identity from explicit projection and factorization hypotheses. \\
    \path{thm:xi:optimal-residual-identity} & partial & needs-strengthening & Lean proves the Loewner-minimal consequence under the elevated algebraic identity; the full identity proof is paper-side. \\
    \path{thm:xi:exact-adequacy} & partial & needs-strengthening & Lean proves the easy implication chain and records hard directions as hypotheses. \\
    \path{thm:xi:adequacy-promotion} & partial & narrowed-surrogate & Lean proves promotion conditional on the upstream adequacy-interface algebraic hypotheses. \\
    \path{thm:xi:chain-rule} & partial & needs-strengthening & Lean proves the Loewner consequence under an elevated Schur chain-rule identity. \\
    \path{cor:xi:same-family-saturation} & partial & narrowed-surrogate & Lean proves saturation under the chain-rule identity and an elevated zero cross-currency hypothesis. \\
    \path{thm:xi:dissolving-post-processing} & theorem & faithful & Lean proves the faithful finite data-processing identity. \\
    \path{thm:xi:native-coarsening} & partial & narrowed-surrogate & Lean proves the Schur-complement difference under the explicit Moore--Penrose comparison hypothesis. \\
    \path{thm:xi:exact-reducing-bridge} & partial & narrowed-surrogate & Lean proves the final residual equality from transported block and Schur-correction identities. \\
    \path{thm:xi:defective-bridge} & partial & narrowed-surrogate & Lean proves the bound from a residual decomposition plus PSD hypotheses for source and defect residuals. \\
    \path{thm:xi:blind-spot-witness} & partial & needs-strengthening & Lean proves the Loewner-violation witness equivalence; the maximum-violation spectral clause is paper-side. \\
    \bottomrule
  \end{tabularx}
  \caption{Manifest coverage at commit \texttt{3bbe9f9}.}
  \label{tab:paper2-lean-coverage}
\end{table}

\begin{table}[H]
  \centering
  \footnotesize
  \setlength{\tabcolsep}{2.5pt}
  \begin{tabularx}{\linewidth}{>{\raggedright\arraybackslash}p{0.38\linewidth} >{\raggedright\arraybackslash}p{0.14\linewidth} >{\raggedright\arraybackslash}X}
    \toprule
    Paper label & Destination & Disclosure note \\
    \midrule
    \path{rmk:xi:null-mode-warning} & appendix & No Lean entry; this standing-hypothesis warning is retained as appendix or setup support. \\
    \path{rmk:xi:projection-interpretation} & appendix & No Lean entry; this interpretive projection remark is handled in prose and may become a short body aside if needed. \\
    \bottomrule
  \end{tabularx}
  \caption{Retained rows with no Lean entry.}
  \label{tab:paper2-nonmechanized-coverage}
\end{table}

\subsection{Lean declaration map}
\label{subsec:paper2-lean-declaration-map}

Table~\ref{tab:paper2-lean-declaration-map} maps each paper label to
its full Lean declaration path; the per-row coverage notes above use
short prose disclosures and refer to this table for the canonical
identifiers.

\begin{table}[H]
  \centering
  \footnotesize
  \setlength{\tabcolsep}{2.5pt}
  \begin{tabularx}{\linewidth}{>{\raggedright\arraybackslash}p{0.34\linewidth} >{\raggedright\arraybackslash}X}
    \toprule
    Paper label & Lean declaration path \\
    \midrule
    \path{thm:xi:minimum-spend-duality} & \path{SixBirdsNeedles.Xi.Currency.xiMinimumSpendDuality} \\
    \path{def:xi:adequacy-residual} & \path{SixBirdsNeedles.Xi.AdequacyResidual.adequacyResidual} \\
    \path{thm:xi:projection-identity} & \path{SixBirdsNeedles.Xi.Projection.projectionIdentity} \\
    \path{thm:xi:optimal-residual-identity} & \path{SixBirdsNeedles.Xi.OptimalResidual.optimalResidualIdentity} \\
    \path{thm:xi:exact-adequacy} & \path{SixBirdsNeedles.Xi.OptimalResidual.exactAdequacy} \\
    \path{thm:xi:adequacy-promotion} & \path{SixBirdsNeedles.Xi.Promotion.adequacyPromotion} \\
    \path{thm:xi:chain-rule} & \path{SixBirdsNeedles.Xi.StrictExtension.chainRule} \\
    \path{cor:xi:same-family-saturation} & \path{SixBirdsNeedles.Xi.StrictExtension.sameFamilySaturation} \\
    \path{thm:xi:dissolving-post-processing} & \path{SixBirdsNeedles.Xi.DataProcessing.dissolvingPostProcessing} \\
    \path{thm:xi:native-coarsening} & \path{SixBirdsNeedles.Xi.DataProcessing.nativeCoarsening} \\
    \path{thm:xi:exact-reducing-bridge} & \path{SixBirdsNeedles.Xi.DataProcessing.exactReducingBridge} \\
    \path{thm:xi:defective-bridge} & \path{SixBirdsNeedles.Xi.DataProcessing.defectiveBridge} \\
    \path{thm:xi:blind-spot-witness} & \path{SixBirdsNeedles.Xi.Obstruction.blindSpotWitness} \\
    \bottomrule
  \end{tabularx}
  \caption{Mapping from paper labels to Lean declaration paths.}
  \label{tab:paper2-lean-declaration-map}
\end{table}
 \section{Predictive Adequacy Ladders}
\label{app:predictive-ladders}

This appendix records the staged propagation behavior of the adequacy residual
when carriers and probe families are indexed by an accepted stage \(j\).  The
forms below parallel the chain-rule and data-processing identities of
\Cref{sec:chain-rule,sec:data-processing}, but they are written as recurrences
rather than as single-step comparison laws.  The body of the paper stays with
the single-step calculus; the ladder forms are retained here for the companion
paper's predictive-membrane framework.

At stage \(j\), write
\[
  \Xi_j=\Xi_{C_j}(D_j\mid L_j),
\]
the stage-local adequacy residual built from the stage carrier
\((E_j,C_j)\), the stage native probe family \(L_j\), and the stage
dissolving probe family \(D_j\).  Each \(\Xi_j\) is therefore the same
Schur-complement residual studied in the body, but computed on the legal
energy quotient at that stage.

The first ladder is multiplicative-additive, in the discrete-Gronwall
tradition of \citet{Pachpatte1998Inequalities}.  Suppose the transported
residuals satisfy
\[
  \Xi_{j+1}\preceqq (1+\eps_j)A_j\Xi_j A_j^*+E_j,
\]
where \(A_j\) is the stage-to-stage transport map and \(E_j\succeq 0\) is
the defect operator paid at stage \(j\).  If the budget recurrence behaves
correspondingly and, after transporting into a common response space,
\[
  \prod_j(1+\eps_j)<\infty,
  \qquad
  \sum_j E_j\preceqq E_\infty,
\]
then predictive adequacy is maintained: one uniform budget bounds the
residual across all accepted stages.  The conjugation by \(A_j\) carries the
residual from one stage to the next, the factor \(1+\eps_j\) records the
allowed multiplicative inflation, and the summability hypotheses say that
neither the inflation nor the accumulated defect can grow without bound.

A stronger ladder assumes strict collapse at each stage:
\[
  \Xi_{j+1}\preceqq (1-\alpha_j)\Xi_j+E_j,
\]
with \(\alpha_j>0\) and
\[
  \prod_j(1-\alpha_j)=0,
\]
together with controlled accumulated defects.  Here \(\alpha_j\) is the
per-stage contraction coefficient.  The vanishing product forces the
multiplicative memory of the initial residual to disappear, while the defect
sum is the remaining obstacle to exact adequacy at the limit.  In the
defect-free case, this is the strict-extension form of blind-spot elimination;
with defects, the limit is controlled modulo the accumulated residual defect.

The multiplicative-additive ladder is the staged analogue of the Schur-update
chain rule in \Cref{thm:xi:chain-rule} together with the data-processing and
transfer laws of \Cref{sec:data-processing}.  The strict-collapse ladder is
the natural staged form of exact adequacy in \Cref{thm:xi:exact-adequacy}: in
the limiting regime, residual disappearance gives the staged ``no leftover''
condition for predictive adequacy.  These ladders are not proved in the body;
they are recorded here because they are used in the predictive-membrane
framework of \citet{Tsiokos2026NeedleKiller}.
 
\clearpage
\bibliographystyle{plainnat}
\bibliography{references}

\end{document}
